% ==========================================================================
%  Section 11 -- The Lean development: organisation and module map.  Owner: W09.
%  Frozen label: sec:lean-overview.  This is the first section of Part II.
%
%  Counts come from generated/metrics.tex only (CONTRACT section 6); no
%  structural count is typed here.  Every machine-checked name below is a
%  declaration of notes/lean_ledger.md section 2 (D12-D14, D30, D32, D42, D48,
%  D57, D58, D60, D61, D63); module names are set in \texttt, because a module
%  is not a declaration and must not be dressed as one.
%
%  Traps observed.  The root aggregator is not claimed to build every module
%  (CONTRACT section 11, final paragraph); the appendices are named in words
%  because their labels live in main.tex, which the paper gate does not scan.
% ==========================================================================
\section{The Lean development: organisation and module map}
\label{sec:lean-overview}

The formalisation is a separate Lake package, \lmod{AImathPPhi2Clean}, whose
sources live under \lmod{formal/lean/s6/pphi2_clean/}, built against Mathlib on
Lean \LeanToolchain{}. \lmod{lake-manifest.json} pins the Mathlib commit
\texttt{85e3a25e}, which is the revision of the tag \lmod{lakefile.lean} names. It has
\NumCleanModules{} modules and
\NumCleanLines{} lines of Lean, together with \NumAuditProbes{} audit probe
files; a comment-stripped count, computed at build time by the paper's own
script and never typed here, finds \NumCleanSorry{} \texttt{sorry} and
\NumCleanProjectAxioms{} project axioms in those modules. The library target
builds: \lmod{lake build AImathPPhi2Clean} elaborates all \NumCleanModules{}
modules with no error on the pinned revision, including the
\NumUnreachableModules{} that the root file does not reach
(\cref{sec:lean-boundary}).  The package sits
inside a larger development: the surrounding directory
\lmod{formal/lean/s6/} holds \NumRootFiles{} further modules with
\NumRootLines{} lines, which this paper does not take as its subject, does not count, and does
not claim to cover.  Every declaration quoted below lives in the namespace
\lmod{Pphi2Clean.Continuum}.

The modules are organised in layers, and the layers follow the order of the
conclusions rather than the order of the proofs: each
one exposes an interface that the next one consumes, so the map below is also a
reading order.  The entry module in the right-hand column is the module to open
first for that layer.

\begin{center}
\scriptsize
\begin{tabular}{@{}p{0.225\linewidth}p{0.325\linewidth}p{0.40\linewidth}@{}}
\toprule
Layer & What it does & Entry module(s) \\
\midrule
(a) periodic torus $L^2$/$H^1$ and the discrete operator &
  Periodic $H^1$ as a closed value-gradient graph completion; the finite
  periodic lattice, its discrete Laplacian with spacing $a$, and an eigenbasis
  of $Q=-\Delta_a+m^2$ &
  \lmod{PeriodicTorusH1}, \lmod{ConcreteLatticeOperator},
  \lmod{PeriodicFourierLaplacian} \\
(b) the massive resolvent and its energy solver &
  The weak massive resolvent from an energy form and a Riesz adjoint; sampled
  and continuum Fourier coefficients; the cutoff tail uniform in the mesh; the
  three covariance limits restated with the resolvent as target &
  \lmod{PeriodicFourierCovarianceLimit},
  \lmod{PeriodicTorusResolventLimit} \\
(c) physical lattice sampling and the Gaussian law &
  Sampling of a test at physical sites with the cell volume as weight; the
  normalised finite covariance; zero mean, characteristic function,
  mixed-moment identity, and the covariance-to-weak-convergence transfer &
  \lmod{PhysicalLatticeGaussianLaw},
  \lmod{PhysicalLatticeSampling} \\
(d) negative-Sobolev carriers, topology and Borel structure &
  The $\HL{2}$ and $\HL{3}$ weights, the completion as $\ell^2$ with its weight
  and deweight maps, the compact-tail comparison, the coordinate embeddings,
  second countability, separability, and the cylinder sigma algebra equal to
  the Borel one &
  \lmod{PeriodicHMinusThreeCompletion},
  \lmod{PeriodicHMinusThreeTopologyBorel} \\
(e) uniform $\HL{2}$ moment and $\HL{3}$ tightness &
  The trace of the squared $\HL{2}$ norm under the finite Gaussian source and
  its mesh-uniform bound; the extended energy, its compact sublevels, the Markov
  tail, and tightness of the law sequence &
  \lmod{PeriodicHMinusTwoUniformTightness},
  \lmod{PeriodicHMinusThreeFixedVolumeTightness} \\
(f) finite real families, characteristic functions and covariance convergence &
  The resolvent covariance matrix, its positive semidefiniteness and Gram
  factor, the canonical target Gaussian, and unconditional weak convergence for
  every finite family of tests &
  \lmod{FiniteRealFamilyWeakLimitInterface},
  \lmod{FiniteFourierCoordinateConvergence} \\
(g) cylinder convergence &
  All the lift maps, the finite-source measure, the paired cosine and sine
  cylinder family, and convergence of every finite Fourier marginal &
  \lmod{PeriodicHMinusThreeCylinderConvergence},
  \lmod{PeriodicHMinusThreeFourierMeasure} \\
(h) Prokhorov subsequence &
  Tightness in the form Prokhorov consumes, and a strictly monotone subsequence
  converging to a Borel probability measure on the carrier &
  \lmod{PeriodicHMinusThreeProkhorovSubsequence} \\
(i) cylinder-to-Borel uniqueness &
  Finite-coordinate cylinders as a $\pi$-system, the Borel-generation step
  discharged from separability, and equality of two Borel probability measures
  with equal indexed cylinder marginals &
  \lmod{PeriodicHMinusThreeCylinderBorelUniqueness} \\
(j) final assembly &
  Inheritance of the target marginals by any subsequential limit, existence and
  uniqueness of the whole-sequence limit, and the packaged
  existence-with-cylinders statement &
  \lmod{PeriodicHMinusThreeWeakLimitMarginal} \\
\bottomrule
\end{tabular}
\end{center}

Layers (a) and (b) are analysis on the torus, layers (c) to (e) are the lattice
laws together with the carrier and the single moment, layers (f) to (i) are the
convergence and identification steps, and layer (j) applies them in order.

The correspondence with Part~I is by result, not by proof step.  Every name in
the right-hand column lives in the namespace \lmod{Pphi2Clean.Continuum} and
the namespace is omitted there.  \Cref{app:ledger} lists the
\NumChainTheorems{} declarations frozen in the two statement contracts, with the
module that declares each of them; the remaining names of the table below are
declarations of the same package whose role is described here rather than
frozen, and their existence is checked against the Lean sources by the paper's
own gate rather than against the registry.

\begin{center}
\footnotesize
\begin{tabular}{@{}p{0.16\linewidth}p{0.82\linewidth}@{}}
\toprule
Result in Part~I & Declaration(s) in \lmod{Pphi2Clean.Continuum} \\
\midrule
\cref{thm:covariance-limit} &
  \lean{fixedVolumeCovarianceInner_tendsto_torusFourierGreenPairing}\newline
  \lean{physicalGJMeasure_fixedVolume_mixedMoment_tendsto_torusFourierGreenPairing}\newline
  \lean{physicalGJMeasure_fixedVolume_covariance_tendsto_torusFourierGreenPairing} \\
\cref{thm:separable} & \lean{periodicHMinusThreeCompletion_separable} \\
\cref{thm:cylinder-borel} &
  \lean{periodicHMinusThree_cylinder_sigma_eq_borel_of_polish} \\
\cref{thm:tightness} &
  \lean{periodicHMinusThree_fixedVolume_freeField_tight} \\
\cref{thm:finite-family} &
  \lean{physicalFamilyEvalProbabilityMeasure_fixedVolume_weak_tendsto_of_periodic_fourier_canonical} \\
\cref{thm:cylinder-uniqueness} &
  \lean{periodicHMinusThree_probabilityMeasure_ext_of_indexed_cylinder_marginals} \\
\cref{lem:subsequence-marginals} &
  \lean{periodicHMinusThree_weak_limit_has_expected_cylinder_marginal}\newline
  \lean{periodicHMinusThree_fixedVolume_subsequential_limit_has_expected_cylinder} \\
\cref{cor:uniqueness} &
  \lean{periodicHMinusThree_fixedVolume_weak_limit_unique} \\
\cref{thm:main} &
  \lean{periodicHMinusThree_fixedVolume_weak_limit_exists_with_expected_cylinders}\newline
  \lean{periodicHMinusThree_fixedVolume_weak_limit_exists}\newline
  \lean{periodicHMinusThree_fixedVolume_weak_limit_has_expected_cylinder} \\
\bottomrule
\end{tabular}
\end{center}

For the \NumChainTheorems{} declarations of the chain the correspondence is
statement by statement: each of the nine results in the left-hand column has at
least one declaration on its right whose conclusion is that statement, and the
three limits inside \cref{thm:covariance-limit} are the three declarations of the
first row, one per line of that theorem.  Two rows carry a qualification, and
they are the two the reader is most likely to overread.
The finite-family declaration consumes the entrywise covariance identification as
a hypothesis --- that is, the very statement \cref{lem:hC-identification} derives
from the covariance limit --- so what the machine closes there is the passage from
that identification to weak convergence of the finite families, with the target
supplied by the Fourier-constructed factor.  The covariance input is exactly
what \cref{thm:finite-family} does not assume, because Part~I spends
\cref{sec:covariance} proving it.
The cylinder-Borel declaration is stated under a Polishness input rather than
unconditionally, and the same module derives that input from separability of the
carrier (\lmod{periodicHMinusThreeCompletion_separable}), which is how the chain
that consumes it discharges it; Part~I's version is unconditional because
\cref{sec:carrier} proves separability before stating it.  Neither qualification
makes a row claim more than Part~I does.
That is the precise
sense in which Part~II mirrors
Part~I, and it is what the table lets a reader check without reading Lean: the
paper's statement on the left, the declaration that carries it in the middle
column, and the module that declares it in Appendix~A.  The row for
\cref{cor:uniqueness} is the weakest entry in the table and must be read as
such: the declaration it names says that two limits of one filter coincide,
which is Hausdorff uniqueness of a filter limit, and not that the limit is
characterised.  The characterisation is the cylinder rows of the table, as it is
in Part~I.  What the table does not
say is anything about the three parts of \cref{thm:main} separately: existence
and the marginal statement are carried by the three declarations of the last row,
while part~(iii), the uniqueness of the limit, is carried by the
cylinder-uniqueness row above it --- that row's declaration is the formal
counterpart of Part~I's part~(iii), though Part~I states no Lean name.

Two qualifications belong here rather than in a footnote, because a reader who
misses either will overread the table.  First, the formal development is a
separate unit from this paper's prose: it has its own module structure, its own
definitions and its own conventions for normalisation, so the package is not the
paper transcribed, and the paper does not claim that the Lean code follows the
proof of Part~I line by line.  Statement identity is what is on offer here, and
it is all that is on offer.  Second, coverage is not uniform across the layers,
and \cref{sec:lean-boundary} collects the five places where the development
diverges from a maximally favourable reading of it: which modules the package
root reaches, a final-theorem-shaped wrapper that assumes its own conclusion,
the total-lift measurability input that is not proved, the two-dimensional
restriction of the measure chain while the covariance theorems are stated for
general dimension, and what a complete build does and does not establish.

Part~II's two appendices are ordered accordingly.  Appendix~A
(\cref{app:ledger}) is the declaration ledger and Appendix~B
(\cref{app:axioms}) the axiom-footprint transcript, and both are generated on
every build from the statement registry and from the recorded transcript rather
than typed, because a hand-maintained list of declaration names is exactly the
kind of object that goes stale without anyone noticing.
